\documentclass[11pt]{article}
\usepackage[margin=1in]{geometry}
\usepackage{graphicx}
\usepackage{booktabs}
\usepackage{hyperref}
\usepackage{xcolor}
\usepackage{amsmath}
\usepackage{amssymb}
\usepackage{caption}
\usepackage{subcaption}
\usepackage{longtable}
\usepackage{placeins}

\title{\textbf{Muon Routing Atlas:}\\Which Transformer Parameters Benefit from Orthogonalized Optimization?}
\author{Pratyaksh Patel}
\date{\today}

\newcommand{\maybeincludegraphics}[2][]{%
  \IfFileExists{#2}{\includegraphics[#1]{#2}}{\fbox{\parbox{0.88\linewidth}{No completed runs were found for this figure at report generation time.}}}%
}
\newcommand{\maybetable}[1]{%
  \IfFileExists{#1}{\input{#1}}{\begin{center}No completed runs were found for this table at report generation time.\end{center}}%
}

\begin{document}
\maketitle

\begin{abstract}
I wanted to understand whether Muon's advantage comes from all transformer matrices or from a smaller subset of modules. I treated the optimizer assignment itself as the experimental variable and built a routing harness around small GPT-style pretraining runs. This report is generated directly from local run artifacts; when an experiment has not been run yet, I state that rather than inventing results.
\end{abstract}

\section{Introduction}

I wanted to understand whether Muon's advantage comes from applying orthogonalized updates to every hidden transformer matrix, or whether most of the gain comes from a smaller subset of modules. The practical version of the question matters: if only feed-forward matrices or value/output attention projections need Muon, then the optimizer can be simpler, cheaper to reason about, and less likely to disturb fragile parts of the model.

The core hypothesis I tested is that Muon may not need to be applied to every hidden matrix. In particular, FFN and V/O matrices may recover much of the gain, while Q/K projections may be less uniformly helpful because they directly shape attention logits and can interact with stability in a different way.

To test this, I treated optimizer routing as the experimental variable. The model, data pipeline, token budget, logging, sampling prompts, and evaluation protocol are kept fixed within each stage, while the set of matrices assigned to Muon changes.

\section{Background}

AdamW is the baseline optimizer for these experiments. It applies adaptive moment estimates with decoupled weight decay and remains a strong default for transformer pretraining.

Muon is an optimizer for hidden matrices that applies momentum and then approximately orthogonalizes the update using a Newton--Schulz iteration \cite{muon_github,muon_writeup}. The motivation is that a matrix update should not concentrate all of its energy in a small number of directions when the parameter itself is a learned linear map.

Routing matters because transformer matrices do not play equivalent roles. Q/K projections determine attention scores, V/O projections move information through attention, and MLP matrices dominate much of the model's nonlinear capacity. Applying the same optimizer to each of these matrices may be unnecessary or even suboptimal.

\section{Method}

\subsection{Model}
I use decoder-only GPT-style transformers with token embeddings, learned positional embeddings, RMSNorm, causal self-attention, and SwiGLU MLPs. Each block exposes explicit parameter names such as \texttt{blocks.3.attn.q\_proj.weight}, \texttt{blocks.3.attn.v\_proj.weight}, and \texttt{blocks.3.mlp.down\_proj.weight}. This naming is part of the experimental design because optimizer assignment is defined over named parameters.

\subsection{Datasets}
TinyStories is used for fast routing ablations \cite{tinystories}. FineWeb-Edu sample-10BT is used for the more serious comparison \cite{fineweb}. FineWeb-Edu is loaded through Hugging Face streaming and can be pre-tokenized into deterministic binary shards \cite{hf_streaming}.

\subsection{Optimizer Routing}
The routing modes are: AdamW everywhere; Muon on all hidden matrices; Muon only on MLP matrices; Muon only on attention matrices; Muon only on V/O matrices; Muon only on Q/K matrices; Muon on everything except Q/K; and layer-restricted early, middle, and late variants.

Embeddings, the language-model head, norms, and biases are always handled by AdamW. Muon is allowed only on selected two-dimensional hidden matrices. Each run writes \texttt{optimizer\_groups.json} before training starts.

\subsection{Training Protocol}
The tokenizer is GPT-2 BPE. The dataloader returns next-token batches where \(x=t_i,\ldots,t_{i+B-1}\) and \(y=t_{i+1},\ldots,t_{i+B}\). I use cosine learning-rate decay with warmup, gradient clipping, mixed precision when available, periodic validation, checkpoints, geometry logging, and deterministic generation prompts.

\subsection{Fair Tuning Protocol}
I kept the tuning budget small and equal across optimizer families. I therefore treat the results as controlled evidence, not as a final statement about the global best hyperparameters.

\maybetable{tables/hyperparameters.tex}

\section{Experiments}

\subsection{Stage 1: TinyStories Routing Sweep}
This stage compares AdamW, full-hidden Muon, MLP-only Muon, attention-only Muon, V/O-only Muon, Q/K-only Muon, and no-QK Muon on TinyStories with two seeds.

\subsection{Stage 2: FineWeb-Edu Serious Comparison}
This stage repeats the most important comparisons on FineWeb-Edu sample-10BT with the mid-sized model. The default token budget is 200M tokens, with support for longer 500M or 1B token runs.

\subsection{Stage 3: Layerwise Atlas}
This stage asks whether early, middle, or late transformer blocks benefit more from Muon, and includes late-only MLP and late-only V/O variants.

\subsection{Stage 4: Batch-Size Stress Test}
This stage compares AdamW, full-hidden Muon, and the best routed Muon configuration across larger effective batch sizes.

\section{Results}

\begin{figure}[ht]
\centering
\maybeincludegraphics[width=0.95\linewidth]{figures/validation_loss_vs_tokens_tinystories.png}
\caption{Validation loss versus tokens on TinyStories.}
\end{figure}

\begin{figure}[ht]
\centering
\maybeincludegraphics[width=0.95\linewidth]{figures/validation_loss_vs_tokens_finewebedu.png}
\caption{Validation loss versus tokens on FineWeb-Edu sample-10BT.}
\end{figure}

\begin{figure}[ht]
\centering
\maybeincludegraphics[width=0.95\linewidth]{figures/validation_loss_vs_wallclock.png}
\caption{Validation loss versus wall-clock time.}
\end{figure}

\begin{figure}[ht]
\centering
\maybeincludegraphics[width=0.95\linewidth]{figures/final_loss_barplot.png}
\caption{Final validation loss by routing mode.}
\end{figure}

\begin{figure}[ht]
\centering
\maybeincludegraphics[width=0.9\linewidth]{figures/routing_heatmap.png}
\caption{Routing heatmap.}
\end{figure}

\begin{figure}[ht]
\centering
\maybeincludegraphics[width=0.9\linewidth]{figures/layerwise_heatmap.png}
\caption{Layerwise routing heatmap.}
\end{figure}

\FloatBarrier

\subsection{Final Results}
\maybetable{tables/final_results.tex}

I summarize the completed runs in the generated tables above. If a stage has not been run yet, the corresponding figure or table says so explicitly.

\section{Update Geometry}

For each routed hidden matrix, I log gradient norm, update norm, weight norm, update-to-weight ratio, top singular value, effective rank, and orthogonality error. Effective rank is computed from normalized singular values \(p_i=s_i/\sum_j s_j\) as
\[
\exp\left(-\sum_i p_i \log p_i\right).
\]
Orthogonality error is computed from \(UU^\top\) or \(U^\top U\), depending on which side is smaller.

\begin{figure}[ht]
\centering
\maybeincludegraphics[width=0.95\linewidth]{figures/update_effective_rank_by_module.png}
\caption{Effective rank by module and routing mode.}
\end{figure}

\begin{figure}[ht]
\centering
\maybeincludegraphics[width=0.95\linewidth]{figures/update_norm_by_module.png}
\caption{Update norm by module and routing mode.}
\end{figure}

\begin{figure}[ht]
\centering
\maybeincludegraphics[width=0.95\linewidth]{figures/singular_value_spectrum_examples.png}
\caption{Example singular-value spectra.}
\end{figure}

\begin{figure}[ht]
\centering
\maybeincludegraphics[width=0.95\linewidth]{figures/qk_logit_growth.png}
\caption{Q/K-related stability diagnostic.}
\end{figure}

\FloatBarrier

The geometry plots are intended to show whether Muon spreads update energy across singular directions, and whether that behavior differs between Q/K, V/O, and MLP matrices.

\section{Qualitative Samples}

At sample intervals and at the final checkpoint, I generate from fixed prompts using temperature \(0.8\), top-\(k=50\), and 120 new tokens. I do not treat these samples as proof of optimizer quality; they are a sanity check for obvious collapse, repetition, or qualitative differences.

\maybetable{tables/qualitative_samples.tex}

If outputs are unavailable, the table states: ``Qualitative samples were not generated for this run.''

\section{Discussion and Limitations}

The main question is whether full-hidden Muon is necessary, or whether no-QK, MLP-only, or V/O-only routing is competitive. If the routed configurations match full-hidden Muon on validation loss and wall-clock efficiency, that would suggest that optimizer design can be more module-aware than a blanket assignment.

The limitations are substantial. The models are small, the experiments are single-GPU, the seed count is limited, and the token budgets are short relative to serious pretraining. TinyStories is synthetic and can reward behavior that does not transfer to broad web text. The hyperparameter budget is deliberately small, so I interpret the results as controlled evidence rather than a global optimizer ranking.

Future versions should test larger models, more seeds, additional architectures such as Mamba or other state-space models, adaptive Muon routing, and more explicit Q/K stabilization.

\section{Conclusion}

I built Muon Routing Atlas to isolate where orthogonalized optimization helps inside GPT-style transformers. The repository trains models, routes Muon across named hidden matrices, logs optimizer groups and update geometry, generates samples, aggregates plots, and compiles this report from local artifacts.

The conclusion should be read from the generated results rather than assumed in advance. If FFN and V/O routing recover most of full-hidden Muon, then Muon may be best understood as a targeted matrix optimizer rather than a uniform replacement for AdamW. If full-hidden Muon remains clearly stronger, then the evidence points toward a broader benefit across transformer hidden matrices.


\bibliographystyle{plain}
\bibliography{references}
\end{document}
